% ============================================================
\chapter{Dualit\'e de Lagrange}
\label{ch:lagrange}
% ============================================================

La dualit\'e lagrangienne est l'une des id\'ees les plus unificatrices de l'optimisation. Son principe~: relaxer les contraintes en les p\'enalisant dans l'objectif via des \emph{multiplicateurs de Lagrange}, puis maximiser sur ces multiplicateurs. Le probl\`eme dual ainsi obtenu est toujours convexe (m\^eme si le primal ne l'est pas), fournit une borne inf\'erieure garantie, et co\"incide avec le primal sous des conditions de \emph{dualit\'e forte} --- lesquelles sont automatiquement satisfaites en optimisation convexe d\`es qu'une condition de qualification de Slater est v\'erifi\'ee. Cette th\'eorie, h\'eriti\`ere des travaux de Lagrange (1797), Kuhn-Tucker (1951) et Rockafellar (1970), est le pont entre la formulation primal et les conditions KKT.

\section{Introduction}

La dualit\'e lagrangienne est un principe fondamental qui associe \`a tout probl\`eme
d'optimisation (primal) un probl\`eme dual. Le dual fournit des bornes inf\'erieures
sur la valeur optimale primale et, sous des conditions de convexit\'e et de qualification,
la dualit\'e forte permet de r\'esoudre le primal via le dual.

\section{Le lagrangien}

\begin{definition}[Lagrangien]
Le \emph{lagrangien} du probl\`eme (P) est la fonction
$L : \R^n \times \R^m_+ \times \R^p \to \R$ d\'efinie par~:
\[
    L(x, \lambda, \nu) = f(x) + \sum_{i=1}^m \lambda_i g_i(x) + \sum_{j=1}^p \nu_j h_j(x).
\]
Les variables $\lambda \in \R^m_+$ et $\nu \in \R^p$ sont les
\emph{variables duales} ou \emph{multiplicateurs de Lagrange}.
\end{definition}

\begin{proposition}[Le lagrangien comme relaxation]
Pour tout $\lambda \geq 0$ et $\nu$, et tout $x$ r\'ealisable~:
\[
    L(x, \lambda, \nu) \leq f(x).
\]
En effet, $\lambda_i g_i(x) \leq 0$ et $\nu_j h_j(x) = 0$.
\end{proposition}

\section{Fonction duale de Lagrange}

\begin{definition}[Fonction duale]
La \emph{fonction duale de Lagrange} est~:
\[
    q(\lambda, \nu) = \inf_{x \in \R^n} L(x, \lambda, \nu)
    = \inf_x \left\{f(x) + \sum_i \lambda_i g_i(x) + \sum_j \nu_j h_j(x)\right\}.
\]
\end{definition}

\begin{theoreme}[Propri\'et\'es de la fonction duale]
\label{thm:duale}
\begin{enumerate}
    \item $q$ est concave (comme infimum de fonctions affines en $(\lambda, \nu)$).
    \item \textbf{Dualit\'e faible}~: $q(\lambda, \nu) \leq p^*$ pour tout
    $\lambda \geq 0$ et $\nu$, o\`u $p^*$ est la valeur optimale primale.
    \item Le domaine $\mathrm{dom}(q) = \{(\lambda,\nu) \mid q(\lambda,\nu) > -\infty\}$
    est convexe.
\end{enumerate}
\end{theoreme}

\begin{proof}[Preuve de la dualit\'e faible]
Pour tout $x$ r\'ealisable et tout $(\lambda, \nu)$ avec $\lambda \geq 0$~:
\[
    q(\lambda, \nu) = \inf_{z} L(z, \lambda, \nu) \leq L(x, \lambda, \nu)
    = f(x) + \sum_i \underbrace{\lambda_i g_i(x)}_{\leq 0}
    + \sum_j \underbrace{\nu_j h_j(x)}_{= 0} \leq f(x).
\]
En prenant l'infimum sur les $x$ r\'ealisables~: $q(\lambda, \nu) \leq p^*$.
\end{proof}

\section{Probl\`eme dual}

\begin{definition}[Probl\`eme dual de Lagrange]
Le \emph{probl\`eme dual} est~:
\[
    (D) \quad d^* = \sup_{\lambda \geq 0, \, \nu} q(\lambda, \nu)
    = \sup_{\lambda \geq 0, \, \nu} \inf_x L(x, \lambda, \nu).
\]
\end{definition}

\begin{definition}[Saut de dualit\'e]
Le \emph{saut de dualit\'e} (\emph{duality gap}) est $p^* - d^* \geq 0$.
Si $p^* = d^*$, on dit que la \emph{dualit\'e forte} a lieu.
\end{definition}

\begin{formules}[title=\textbf{Relation primal--dual}]
\[
    d^* = \sup_{\lambda \geq 0, \nu} \inf_x L(x,\lambda,\nu)
    \leq \inf_x \sup_{\lambda \geq 0, \nu} L(x,\lambda,\nu) = p^*.
\]
L'in\'egalit\'e est l'in\'egalit\'e max-min.
\end{formules}

\section{Dualit\'e forte}

\begin{theoreme}[Dualit\'e forte --- Slater]
\label{thm:slater}
Si le probl\`eme primal est convexe ($f, g_i$ convexes, $h_j$ affines), que $p^*$
est fini, et que la condition de Slater est satisfaite, alors~:
\begin{enumerate}
    \item La dualit\'e forte a lieu~: $p^* = d^*$.
    \item Le supremum dual est atteint~: il existe $(\lambda^*, \nu^*)$ optimal.
\end{enumerate}
\end{theoreme}

\begin{proof}[Id\'ee de la preuve]
On consid\`ere l'ensemble~:
\[
    \mathcal{G} = \{(u, v, t) \in \R^m \times \R^p \times \R \mid
    \exists x : g_i(x) \leq u_i, h_j(x) = v_j, f(x) \leq t\}.
\]
$\mathcal{G}$ est convexe (convexit\'e de $f, g_i$ et affinit\'e de $h_j$).
Le point $(0, 0, p^*)$ est sur le bord de $\mathcal{G}$. Par le th\'eor\`eme de
l'hyperplan d'appui, il existe un hyperplan s\'eparant, ce qui donne les
multiplicateurs de Lagrange optimaux.
\end{proof}

\section{Exemples de calcul dual}

\begin{exemple}[Programmation lin\'eaire]
Primal~: $\min c^Tx$ sous $Ax \leq b$, $x \geq 0$.

$L(x, \lambda, \mu) = c^Tx + \lambda^T(Ax - b) - \mu^Tx = (c + A^T\lambda - \mu)^Tx - \lambda^Tb$.

$q(\lambda, \mu) = \inf_x L = \begin{cases} -\lambda^Tb & \text{si } c + A^T\lambda - \mu = 0, \\ -\infty & \text{sinon.}\end{cases}$

Dual~: $\max -b^T\lambda$ sous $A^T\lambda \leq c$, $\lambda \geq 0$,
soit $\max b^T y$ sous $A^T y \leq c$, $y \leq 0$ (en posant $y = -\lambda$).

Apr\`es reformulation standard~: le dual de la PL est une PL.
\end{exemple}

\begin{exemple}[Norme minimale]
Primal~: $\min \norm{x}^2$ sous $Ax = b$.

$L(x, \nu) = \norm{x}^2 + \nu^T(Ax - b)$.

$\inf_x L = \inf_x \{\norm{x}^2 + \nu^TAx\} - \nu^Tb$.

La condition d'optimalit\'e $2x + A^T\nu = 0$ donne $x = -\frac{1}{2}A^T\nu$.

$q(\nu) = \frac{1}{4}\nu^TAA^T\nu - \frac{1}{2}\nu^TAA^T\nu - \nu^Tb = -\frac{1}{4}\nu^TAA^T\nu - \nu^Tb$.

Dual~: $\max -\frac{1}{4}\nu^TAA^T\nu - \nu^Tb$.
\end{exemple}

\begin{exemple}[Entropie maximale]
$\min \sum_i x_i \log x_i$ sous $Ax = b$, $\mathbf{1}^Tx = 1$, $x \geq 0$.

Dual~: $\max -\log\left(\sum_i \exp(-a_i^T\nu - \mu)\right)$ sous les contraintes
duales, o\`u $a_i$ est la $i$-\`eme colonne de $A$.
\end{exemple}

\section{Interpr\'etation min-max}

\begin{theoreme}[Th\'eor\`eme du minimax de Von Neumann--Sion]
Si $\mathcal{X}$ est compact convexe, $\mathcal{Y}$ est convexe, et $\Phi(x,y)$
est convexe-concave (convexe en $x$, concave en $y$), alors~:
\[
    \min_{x \in \mathcal{X}} \sup_{y \in \mathcal{Y}} \Phi(x,y)
    = \sup_{y \in \mathcal{Y}} \min_{x \in \mathcal{X}} \Phi(x,y).
\]
\end{theoreme}

\begin{remarque}
Le lagrangien $L(x, \lambda, \nu)$ est convexe en $x$ (somme de fonctions convexes)
et affine (donc concave) en $(\lambda, \nu)$. Sous les conditions de Slater, le
minimax s'applique, donnant la dualit\'e forte.
\end{remarque}

\section{Point-selle du lagrangien}

\begin{definition}[Point-selle]
$(x^*, \lambda^*, \nu^*)$ est un \emph{point-selle} du lagrangien si~:
\[
    L(x^*, \lambda, \nu) \leq L(x^*, \lambda^*, \nu^*) \leq L(x, \lambda^*, \nu^*),
    \quad \forall x, \; \forall \lambda \geq 0, \nu.
\]
\end{definition}

\begin{theoreme}[\'Equivalence point-selle et KKT]
Sous les conditions de Slater, $(x^*, \lambda^*, \nu^*)$ est un point-selle du
lagrangien si et seulement si $x^*$ est primal optimal, $(\lambda^*, \nu^*)$ est
dual optimal, et la dualit\'e forte a lieu.
\end{theoreme}

\section{Interpr\'etation \'economique}

\begin{intuition}
Consid\'erons le probl\`eme de maximisation de profit d'une entreprise sous contraintes
de ressources~:
\[
    \max_x \; \mathrm{profit}(x) \quad \text{sous} \quad
    \mathrm{ressource}_i(x) \leq b_i.
\]
Le multiplicateur $\lambda_i^*$ repr\'esente le \emph{prix implicite} de la
ressource $i$~: c'est le gain marginal obtenu en augmentant la disponibilit\'e
de la ressource $i$ d'une unit\'e. Un march\'e concurrentiel fixerait naturellement
ce prix.
\end{intuition}

\section{Impl\'ementation Python}

\begin{codeblock}[title=\textbf{Dualit\'e lagrangienne et saut de dualit\'e}]
\begin{minted}{python}
import numpy as np
import cvxpy as cp

# Exemple: min ||x||^2 s.t. Ax = b
np.random.seed(42)
m, n = 3, 10
A = np.random.randn(m, n)
b = np.random.randn(m)

# Primal
x = cp.Variable(n)
primal = cp.Problem(cp.Minimize(cp.sum_squares(x)), [A @ x == b])
primal.solve()
p_star = primal.value
print(f"Valeur primale p* = {p_star:.6f}")
print(f"Multiplicateurs: {primal.constraints[0].dual_value}")

# Dual analytique: max -1/4 nu^T A A^T nu - nu^T b
nu = cp.Variable(m)
M = A @ A.T
dual = cp.Problem(cp.Maximize(-0.25 * cp.quad_form(nu, M) - b @ nu))
dual.solve()
d_star = dual.value
print(f"Valeur duale d*   = {d_star:.6f}")
print(f"Saut de dualite   = {p_star - d_star:.2e}")

# Programmation lineaire: verification dualite forte
c = np.random.randn(n)
A_ineq = np.random.randn(m, n)
b_ineq = np.abs(np.random.randn(m)) + 1

x_lp = cp.Variable(n)
primal_lp = cp.Problem(cp.Minimize(c @ x_lp),
                        [A_ineq @ x_lp <= b_ineq, x_lp >= 0])
primal_lp.solve()

y_lp = cp.Variable(m)
dual_lp = cp.Problem(cp.Maximize(b_ineq @ y_lp),
                      [A_ineq.T @ y_lp + c >= 0, y_lp <= 0])
dual_lp.solve()

print(f"\nPL primal: {primal_lp.value:.6f}")
print(f"PL dual:   {dual_lp.value:.6f}")
print(f"Gap:       {abs(primal_lp.value - dual_lp.value):.2e}")
\end{minted}
\end{codeblock}

\section{Exercices}

\begin{exercice}[$\star$]
\'Ecrire le dual de Lagrange de $\min c^Tx$ sous $Ax = b$, $x \geq 0$.
\end{exercice}

\begin{exercice}[$\star$]
Calculer la fonction duale pour $\min \frac{1}{2}x^TQx + c^Tx$ sous $Ax = b$
avec $Q \succ 0$.
\end{exercice}

\begin{exercice}[$\star\star$]
Montrer que la dualit\'e forte a lieu pour la programmation lin\'eaire sans
n\'ecessiter la condition de Slater.
\end{exercice}

\begin{exercice}[$\star\star$]
D\'eriver le dual du probl\`eme SVM \`a marge souple et montrer qu'il s'agit d'un
QP born\'e.
\end{exercice}

\begin{exercice}[$\star\star$]
Montrer que $(x^*, \lambda^*, \nu^*)$ est un point-selle du lagrangien si et seulement
si les conditions KKT sont satisfaites.
\end{exercice}

\begin{exercice}[$\star\star\star$]
D\'emontrer le th\'eor\`eme de Slater en utilisant le th\'eor\`eme de l'hyperplan
d'appui appliqu\'e \`a l'ensemble $\mathcal{G}$ d\'efini dans la preuve.
\end{exercice}

\begin{exercice}[$\star\star\star$]
Soit $p^*(u)$ la valeur optimale du probl\`eme perturb\'e
$\min f(x)$ sous $g_i(x) \leq u_i$. Montrer que $p^*$ est convexe en $u$.
En d\'eduire l'interpr\'etation des multiplicateurs comme sous-gradients de $p^*$.
\end{exercice}
